\documentclass[11pt]{article}
\usepackage[a4paper,margin=2.2cm]{geometry}
\usepackage[T1]{fontenc}
\usepackage{textcomp}
\usepackage[spanish]{babel}
\usepackage{amsmath,amssymb}
\usepackage{enumitem}
\usepackage{tikz}
\usetikzlibrary{calc,arrows.meta,patterns,decorations.pathmorphing}
\usepackage{xcolor}
\usepackage{colortbl}
\usepackage{multicol}
\usepackage{parskip}
\setlength{\parskip}{6pt}
\usepackage{lmodern}
\renewcommand{\familydefault}{\sfdefault}

\definecolor{unalblue}{RGB}{31,73,124}

\setlist[enumerate,1]{itemsep=10pt,topsep=8pt,leftmargin=2.2em,label=\protect\fbox{\arabic*}}
\newlist{opciones}{enumerate}{1}
\setlist[opciones]{label=(\alph*),itemsep=8pt,topsep=6pt,leftmargin=2em,font=\normalfont}

\newcommand{\ptsbox}[1]{\fboxsep=3pt\fbox{\footnotesize #1}\hspace{4pt}}
\newcommand{\borradorbox}[1][3cm]{%
  \noindent\fbox{\begin{minipage}[t][#1][t]{0.97\linewidth}
  {\scriptsize\itshape Espacio borrador, nada escrito aquí será calificado.}
  \end{minipage}}
}

\newcommand{\vect}[2]{\begin{pmatrix}#1\\#2\end{pmatrix}}
\newcommand{\vectiii}[3]{\begin{pmatrix}#1\\#2\\#3\end{pmatrix}}

\newcommand{\examheader}{%
  \noindent
  \begin{minipage}[t]{0.6\linewidth}
    {\small\bfseries Geometría Vectorial --- Parcial 2}\\
    {\small Octubre de 2022}
  \end{minipage}%
  \hfill
  \begin{minipage}[t]{0.37\linewidth}
    \raggedleft
    \fbox{\begin{minipage}{0.95\linewidth}\centering\scriptsize Escuela de Matemáticas. Universidad Nacional de Colombia, Sede Medellín.\end{minipage}}
  \end{minipage}\\[4pt]
  \rule{\linewidth}{0.6pt}
}
\newcommand{\examfooter}{%
  \vfill
  \rule{\linewidth}{0.4pt}\\[1pt]
  {\scriptsize\textit{Transcripción digital --- MathUNAL}\hfill\textcolor{unalblue}{\textbf{\thepage}}}
}
\pagestyle{empty}

\begin{document}
\examheader

\begin{center}
{\bfseries Geometría Vectorial --- Parcial 2}\\
Octubre de 2022

\vspace{6pt}
\textbf{Instrucciones}
\end{center}

La duración del examen es \textbf{1:50 horas}. Verifique que sus datos de identificación son correctos. La prueba consta de cinco preguntas. Escriba sus respuestas a los ejercicios 1,2,3,4 en la tabla que está abajo en esta página. Solo las respuestas anotadas en la tabla serán tenidas en cuenta. No use fracciones, solamente decimales y redondee a dos cifras decimales, por ejemplo, si su respuesta es $\frac{1}{3}$ escriba 0.33. La pregunta 5 es de procedimiento, y su respuesta debe escribirse debajo del enunciado. \textbf{No} está permitido: hacer preguntas, usar calculadoras ni trabajar en hojas diferentes a las del examen. Por favor no desengrapar el examen y entregarlo completo al final del examen al vigilante. No se permite sacar el teléfono celular durante el examen.

\vspace{8pt}
\begin{center}
\textbf{\textsc{Tabla puntaje y respuestas}}

\vspace{4pt}
\begin{tabular}{|c|c|l|}
\hline
Ejercicio & Porcentaje & Respuesta\\
\hline
1 & 17\% & a\quad b\quad c\quad d\quad e\quad f \\[6pt]
\hline
2 & 18\% & \\[14pt]
\hline
3 & 17\% & \\[14pt]
\hline
4 & 18\% & \\[14pt]
\hline
5 & 30\% & \cellcolor{black!15}Procedimiento \\[6pt]
\hline
\end{tabular}
\end{center}

\vspace{10pt}

\begin{enumerate}

\item Una recta pasa por el punto $P = \vect{10^5-1}{10^5-4}$ y tiene pendiente $4$. Señalar de los siguientes puntos, aquel único que también pertenece a la recta:

\begin{opciones}
\item $\vect{10^5}{10^5}$
\item $\vect{10^5-4}{10^5+1}$
\item $\vect{10^5+1}{10^5-4+2}$
\item $\vect{1}{4}$
\item $\vect{10^5-1+4}{10^5-4+1}$
\item Ninguno de los anteriores.
\end{opciones}

\vspace{40pt}

\item Sea $X = \vect{-1}{-5}$. Suponga que $\mathcal{L}$ es una recta con vector director $D$ y vector normal $N$, y que $Q$ es un punto en $\mathcal{L}$. Si $P_D(X) = \vect{-3}{-3}$, $P_N(Q) = \vect{1}{-1}$, determine $P_N(X)$.

\newpage
\examheader

\item Considere la circunferencia $\mathcal{C}$ con centro en $\vect{-1}{2}$ y radio 5 y la recta $\mathcal{L}$ paralela al eje $x$ y con intercepto 6. Hallar la longitud de la cuerda comprendida entre los puntos de intersección de $\mathcal{C}$ con $\mathcal{L}$.

\vspace{100pt}

\item Considere la transformación lineal
\[
T\vect{x}{y} = \vect{16x+(-4)y}{-4x+y}.
\]
El núcleo de $T$ es la recta que pasa por el origen y tiene vector director $\vect{1}{k}$. Determine el valor de $k$.

\newpage
\examheader

\item \textbf{Pregunta de procedimiento:} En el espacio de abajo escriba un procedimiento completo justificando cada paso. Sea $T$ la transformación lineal definida por
\[
T\vect{x}{y} = \vect{2x-y}{x+y},
\]
y sea $\mathcal{L}$ la línea recta con ecuación general $3x-y-2=0$. Hallar una ecuación general de la línea recta $\mathcal{L}' = T(\mathcal{L})$.

\vspace{200pt}

\end{enumerate}

\examfooter
\end{document}
